\subsection{What an integral Witt coefficient does not detect}

The following is a failure gate for a proposed same-discriminant selection
argument.  It neither proves nor disproves existence of a different
auxiliary field with all the required units.

\begin{proposition}[KT-WITT-PULLBACK-MASS-OBSTRUCTION]
Let $F_{31}$ be the normalized degree-six Eisenstein series of
KT-SENARY-GENUS-TRACE-DETECTOR, so $b=17$ and $k=20$.  Put
\[
 C=\begin{pmatrix}2&-1&0\\-1&2&-1\\0&-1&2\end{pmatrix},
 \qquad
 B=\begin{pmatrix}2&-1&0\\-1&2&-1\\0&-1&14\end{pmatrix}.
\]
For the Witt restriction
$W(F)(Z_1,Z_2)=F(\operatorname{diag}(Z_1,Z_2))$, define the
\emph{determinant-$79$ contribution} to its $(C/2,B/2)$ coefficient by
\[
 c_{79}=
 \sum_{\substack{R\in M_3(\mathbb Z)\\
 G_R=\left(\begin{smallmatrix}C&R\\R^t&B\end{smallmatrix}\right)>0\\
 \det G_R=79}}
 a(F_{31},G_R/2).
 \tag{WP.1}
\]
There are exactly $56$ summands, and
\[
 c_{79}\equiv27\pmod{31},\qquad
 \mathcal T_{-79}(F_{31})\equiv0\pmod{31}.
 \tag{WP.2}
Thus even a unit determinant-filtered contribution to an ordinary Witt
coefficient does not force the product
$\beta_{17}(\chi_{-79})\beta_3(\chi_{-79})$ to be a unit.
The determinant filter in (WP.1) is just a finite coefficient selection;
no modularity of this filtering operation is asserted.
\end{proposition}

\begin{proof}
The coefficient formula for a Witt restriction follows immediately by
substituting a block diagonal argument in the Fourier expansion.  It sums
over the cross-block matrices $R$, including positive semidefinite matrices
in the unfiltered coefficient.  For the strictly positive matrices in
(WP.1), the two-by-two principal minors give
$|R_{ij}|^2<C_{ii}B_{jj}$.  The first two columns therefore have entries
in $\{-1,0,1\}$ and the third has entries in $\{-5,\ldots,5\}$.
This finite box has $3^6 11^3=970299$ elements.

Here $\det C=4$, $\det B=40$, and
\[
 M=4C^{-1}=\begin{pmatrix}3&2&1\\2&4&2\\1&2&3\end{pmatrix}.
\]
Writing the columns of $R$ as $x,y,v$, positivity requires
$x^tMx<8$, $y^tMy<8$, and $v^tMv<56$.  These leave respectively
$15,15,411$ vectors.  The scaled Schur complement is
\[
 S=4B-R^tMR=\begin{pmatrix}a&z&s\\z&b&t\\s&t&c\end{pmatrix},
 \qquad \det G_R=\det(S)/16.
\]
Sylvester's criterion says that $G_R>0$ exactly when
$a>0$, $ab-z^2>0$, and $\det S>0$.  The saved exact certificate enumerates
these vectors and checks
\[
 (ab-z^2)c-at^2-bs^2+2zst=16\cdot79.
\]
It obtains $56$ solutions.  A separate enumeration of the full coordinate
box uses integer Bareiss determinants of the four-, five-, and six-dimensional
principal matrices and obtains the identical list.  Thus both the finiteness
and completeness of the count are explicit.

Since $-79$ is fundamental, the coefficient formula already used in
KT-SENARY-GENUS-TRACE-DETECTOR gives
$a(F_{31},G_R/2)=-\beta_{17}(\chi_{-79})$ for every summand, independently
of its integral isometry class.  The previously proved exact values are
$\beta_{17}(\chi_{-79})\equiv20\pmod{31}$ and
$\beta_3(\chi_{-79})=992$; the new certificate also recomputes both
from the rational Bernoulli-polynomial formula.  Hence
$c_{79}\equiv-56\cdot20=27\pmod{31}$.
The genus trace is zero because its mass is $31/720$, by (SG.3).
\end{proof}

\begin{lemma}[KT-GENUS-PROJECTOR-INTEGRALITY]
Let $p\geq31$ and let $\mathcal L_D$ be a senary genus as in
KT-SENARY-GENUS-TRACE-DETECTOR.  On
$\Lambda=\bigoplus_{[L]}\mathbb Z_{(p)}\delta_L$, consider the canonical
mass-orthogonal projector onto constant functions:
\[
 P(f)=\frac{\sum_{[L]}f(L)/|\operatorname O(L)|}
 {\sum_{[L]}1/|\operatorname O(L)|}\,\mathbf1.
 \tag{WP.3}
\]
Then $P(\Lambda)\subseteq\Lambda$ if and only if
$\beta_3(\chi_D)$ is a $p$-unit.
\end{lemma}
\begin{proof}
All automorphism denominators are $p$-units, as proved in the cited
senary result.  Therefore the mass $M$ is $p$-integral and every numerator
$\langle\delta_L,\mathbf1\rangle=1/|\operatorname O(L)|$ is a unit.
Formula (WP.3) preserves the lattice exactly when $M$ is a unit.
Equation (SG.3) identifies this condition with the asserted Bernoulli unit.
For $p=31,D=-79$, its value on every $\delta_L$ has valuation $-1$.
This concerns the canonical mass-orthogonal projector; arbitrary
nonorthogonal projections onto the constant line do exist integrally.
\end{proof}

\paragraph{The remaining pullback requirement.}
An unnormalized integral coefficient functional could still prove a useful
simultaneous-unit result.  It must actually equal an integral combination
of the genus traces, and its unit value must be established.  Ordinary
cross-block summation weights lattices by their numbers of compatible
ordered bases; (WP.2) shows that these weights cannot simply be replaced
by the mass weights.  Normalizing a genus average by its total mass, as in
(WP.3), already requires the third unit.  Neither observation is counted
as movement on the remaining arithmetic gate.

\paragraph{Published integrality scope.}
Pitale--Saha--Schmidt, \emph{Integrality and cuspidality of pullbacks of
nearly holomorphic Siegel Eisenstein series}, Publ. Mat. \textbf{66}
(2022), 405--434, DOI 10.5565/PUBLMAT6612216, Theorem 1.1(ii) on p.~408
and Remark 1.2 on p.~409, give a general normalized coefficient-integrality
theorem with hypothesis $p\geq2k$.  It cannot be invoked directly at
$k=(p+9)/2$.  This is an applicability statement about that theorem,
not a claim that our holomorphic coefficients fail to be integral; the
coefficients in (WP.1) are integral by their explicit formula.  The
theorem statement and explanatory remark were read for this scope check;
no fresh review of its full proof is claimed.
